\section{Introduction}
\label{sec:introduction}

An LLM request is not a single GPU operation. It begins with a parallel
prefill over the prompt and continues as an autoregressive sequence of decode
iterations. Serving engines exploit this structure by admitting and retiring
requests between iterations. Orca introduced iteration-level scheduling
\citep{yu2022orca}; \vllm{} couples continuous scheduling with PagedAttention
to reduce key--value-cache waste \citep{kwon2023pagedattention}. These engines
therefore already decide which requests share the GPU at fine granularity.

The deployment stack often adds a second scheduler. An HTTP gateway may hold
compatible requests for a few milliseconds, concatenate their prompts into one
backend call, and then split the outputs. This pattern is attractive because
batching can amortize HTTP, tokenization, and launch overhead, and it has a long
history in general prediction serving \citep{crankshaw2017clipper}. But in
front of a continuously batching LLM engine, the gateway does not create GPU
batching from nothing. It changes the time and shape in which work becomes
visible to an engine that already batches. We call this interaction
\emph{double batching}.

Double batching creates a non-obvious trade-off. Waiting can expose several
prompts in one HTTP request, but it also withholds those prompts from the
backend scheduler. A longer gateway batch can reduce per-prompt backend
overhead while worsening time to first token and head-of-line blocking. The
gateway sees only request-level completion, whereas the engine sees prefill,
decode, KV-cache pressure, and token-level progress. Optimizing the outer
scheduler without measuring the inner scheduler can therefore optimize the
wrong bottleneck.

This paper tests one falsifiable systems hypothesis:

\begin{quote}
\textbf{H:} At offered loads that direct \vllm{} can fully admit, a millisecond
gateway window is too short relative to inter-arrival time to form useful
batches. Multi-request outer batches, when they appear, are primarily backlog
coalescing caused by the gateway awaiting a backend call. Consequently, outer
batching improves latency only if the backend-time reduction from coalescing
exceeds the queueing delay and loss of scheduling freedom imposed on \vllm{}.
\end{quote}

The null hypothesis is operationally useful: fixed and adaptive gateway
batching provide no repeatable tail-latency advantage over pass-through at
matched offered load. A positive result requires a paired, run-level
improvement visible across repetitions and a latency decomposition that repays
added queueing; a negative result is evidence that the engine, not the gateway,
should own scheduling for the tested regime.

We study the question with \adip{}, an asynchronous gateway and measurement
artifact. The final matrix compares four paths---direct \vllm{}, gateway
pass-through, fixed-window batching, and a frozen adaptive policy---at
0.5, 1.0, 1.5, and 1.8 requests/s. Every condition uses identical prompts,
generation parameters, open-loop arrivals, and a single consumer-grade GPU.
Three seeds serve as the replication unit. We retain request-level latency,
queue and backend decomposition, batch identifiers and sizes, available GPU
telemetry, failures, and arrival drift.

The paper makes four contributions:

\begin{enumerate}[leftmargin=1.4em]
  \item It identifies a precise gap between engine-level continuous batching
        and request-level gateway batching, and gives a simple model separating
        \emph{window formation} from \emph{backlog formation}.
  \item It provides a tested implementation of fixed and adaptive outer
        scheduling with compatibility keys, deadline-aware closure, structured
        telemetry, and explicit cancellation and shutdown semantics.
  \item It evaluates the schedulers using predeclared open-loop arrivals and
        run-level uncertainty rather than treating requests from one run as
        independent replications.
  \item It records the complete validity trail, including interrupted and
        harness-rejected attempts, and binds final claims to machine-readable
        analysis artifacts so neutral or negative findings remain auditable.
\end{enumerate}

The scope is intentionally narrow. We do not propose a replacement for
\vllm{}, claim production-scale generality, or tune on the final matrix. The
goal is a defensible measurement study that explains a common layered-serving
decision on hardware available to an individual researcher.
